\begin{abstract}
When an LLM writes code to satisfy a multi-step specification, does it know \emph{what it intends to do next} before it starts writing? We show that code LLMs encode a compact representation of their next instrumental sub-goal before emitting its implementation, and that this encoding is \emph{decodable} and \emph{portable}. Using causal mediation analysis, we localize this encoding to a set of attention heads, which we call the \goalheads. We introduce \goalscope, a lightweight probe that uses what these heads write to the residual stream to verbalize the upcoming sub-task in natural language. We demonstrate that the encoding is abstract: it can be extracted from one context and transplanted to trigger the same sub-goal in a different context and programming language. These heads also mediate the model's own implementation: preserving their access to task information largely preserves the capability of executing the next sub-task, while restricting that access in a broader set of goal heads sharply impairs it. We replicate the readout and transfer findings across model families and sizes, and explore the same heads' readouts in non-code settings. Our results provide concrete, testable evidence that LLMs encode abstract representations of intermediate goals before realizing them in tokens.
\end{abstract}
